
\paragraph{Impossibility of Satisfying Multiple Equal-Cost Constraints.}
\label{ssec:impossibility}
It is natural to argue there might be multiple cost functions that we
would like to equalize across groups.
However, satisfying more than one distinct
equal-cost constraint (i.e. different curves in the F.P./F.N. plane) is infeasible.
%However,
%the single-cost-constraint relaxation is the strongest feasible notion of
%calibrated non-discrimination (in the sense that multiple distinct equal-cost
%constraints cannot be simultaneously satisfied).
%
\begin{thm}[Generalized impossibility result] \label{thm:impossibility}
  Let $h_1$ and $h_2$ be calibrated classifiers for $G_1$ and $G_2$ with equal
  cost with respect to
  $\fairconstname$.
  If $\mu_1 \ne \mu_2$, and if $h_1$ and $h_2$ also have equal cost with
  respect to a different cost function $\fairconstname'$, then $h_1$ and $h_2$
  must be perfect classifiers.
\end{thm}
%
(Proof in \autoref{app:impossiblility}). Note that this is
a generalization of the impossibility result of \cite{kleinberg2016inherent}.
Furthermore, we show in \autoref{thm:approx-impossibility} (in \autoref{app:impossiblility})
that this holds in an
approximate sense: if calibration and multiple distinct equal-cost constraints
are approximately achieved by some classifier, then that classifier must have
approximately zero generalized false-positive and false-negative rates.


%If, instead, we forego the calibration condition in favor of maintaining both
%the false-positive and false-negative conditions, we obtain the Equalized Odds framework
%of \citet{hardt2016equality}.
%%
%Thus our framework in combination with Equalized Odds encompass all possible relaxations
%of the \citeauthor{kleinberg2016inherent} conditions.
%
%While our definition of fairness may not be appropriate for some classification
%tasks, this framework may prove useful in settings where each error type is undesirable, or where
%we care about matching only one specific error type.
%In settings in which both the false-positive and false-negative conditions are absolutely necessary,
%then the Equalized Odds framework may be more applicable.

